% ============================================================================
\section{Introducción}
% ============================================================================

\subsection{Qué pide la práctica}

El enunciado plantea el \emph{problema de cubrimiento}
\[
  \min\ \{\, z = c^{\top}x \;\mid\; Ax \ge \mathbf{1},\ x \text{ binaria} \,\}
\]
con una matriz $A$ de $m$ filas y $n$ columnas y un vector de costes $c$, cuyos datos están
en los ficheros \cod{SCPXX.dat}, y pide implementar:
\begin{enumerate}
  \item un algoritmo \textbf{constructivo} y uno de \textbf{búsqueda local};
  \item las metaheurísticas \textbf{Temple Simulado} y \textbf{Búsqueda Tabú};
  \item las metaheurísticas \textbf{GRASP} y \textbf{Algoritmo Genético};
  \item una metaheurística de \textbf{Búsqueda en Entorno Variable} (VNS);
  \item un informe con los algoritmos, el código, la \textbf{calibración de los parámetros},
        las pruebas con las instancias y el análisis de los resultados.
\end{enumerate}

Esta memoria cubre los algoritmos, el código y la calibración (incluida la validación de la
calibración). El informe final de resultados y la presentación se apoyarán en ella.

\subsection{Qué se entrega}

\begin{center}
\begin{tabular}{@{}lp{11.2cm}@{}}
\toprule
Carpeta & Contenido \\ \midrule
\cod{codigo/} & Programa en C (diez ficheros), comentado línea a línea, y el guion
                \cod{compilar.sh} para compilarlo. Un único ejecutable, \cod{scp}, contiene
                los siete algoritmos. \\
\cod{experimentos/} & Guiones en Python que \emph{solo} lanzan el ejecutable muchas veces y
                recogen lo que devuelve (los algoritmos están íntegramente en C), junto con
                todos los ficheros de resultados. \\
\cod{memoria/} & Este documento en \LaTeX{} y las tablas y figuras que se generan a partir de
                los resultados. \\
\cod{datos/} & Las 6 instancias usadas para \textbf{calibrar}. \\
\cod{Validacion\_datos/} & Las 24 instancias usadas para \textbf{validar}. \\
\bottomrule
\end{tabular}
\end{center}

\subsection{Cómo leer este documento}

El documento está escrito para alguien que no conoce de antemano ni las metaheurísticas ni el
lenguaje C. Por eso sigue siempre el mismo orden:
\begin{itemize}
  \item La sección~\ref{sec:problema} explica el problema y las instancias.
  \item La sección~\ref{sec:conceptos} da el vocabulario mínimo (solución, entorno, óptimo
        local, intensificar y diversificar) que se usa en todo lo demás.
  \item La sección~\ref{sec:programa} describe las piezas del programa que comparten todos
        los algoritmos. Conviene leerla antes que las siguientes.
  \item Las secciones~\ref{sec:constructivo} a~\ref{sec:vns} dedican un apartado a cada
        algoritmo con la misma estructura: \emph{idea}, \emph{esquema de los apuntes},
        \emph{adaptación al problema} (con la justificación de cada decisión),
        \emph{parámetros} y \emph{recorrido por el código}.
  \item Las secciones~\ref{sec:metodologia} a~\ref{sec:validacion} explican cómo se han
        elegido los valores de los parámetros y cómo se ha comprobado que la elección es
        buena. Ahí se responde con detalle a la pregunta de \emph{qué instancias se usan para
        calibrar y cuáles para validar}.
  \item El apéndice contiene el código completo.
\end{itemize}

Tres tipos de recuadro ayudan a localizar lo importante:

\begin{idea}
Resume en dos o tres frases la idea que hay que retener.
\end{idea}
\begin{apuntes}
Reproduce el esquema del algoritmo tal como aparece en los apuntes de la asignatura. El
código sigue ese esquema paso a paso.
\end{apuntes}
\begin{decision}
Explica una decisión que los apuntes dejan abierta (porque depende del problema) y por qué se
ha resuelto de esa manera.
\end{decision}

\paragraph{Convenio numérico.} Los decimales se escriben con punto (4.13 y no 4,13) porque
así los escribe el programa y así aparecen en los ficheros de resultados.
